\chapter{算法设计与实验}

\section{问题建模}

以连续函数优化问题为例，设搜索空间为 $\Omega\subseteq\mathbb{R}^{d}$，算法目标是在给定迭代次数内寻找使目标函数取得较小值的候选解。算法研究类论文应明确问题输入、输出、约束条件、终止条件和评价指标。图\ref{fig:framework}给出了一个典型算法研究流程示例。

\begin{figure}[htbp]
  \centering
  \fbox{\parbox[c][4.2cm][c]{0.82\textwidth}{\centering 算法研究流程图占位\\问题建模、算法设计、参数设置、实验比较}}
  \caption{算法研究流程示意图}
  \label{fig:framework}
\end{figure}

\section{改进策略}

粒子群优化算法通过模拟群体协作搜索过程更新粒子位置。基础算法通常包含种群初始化、适应度计算、个体最优更新、全局最优更新和速度位置更新等步骤。为了提升算法在复杂多峰函数上的搜索能力，可引入动态惯性权重和局部扰动策略，在保持全局探索能力的同时增强后期局部开发能力。相关综述通常会从参数控制、拓扑结构和混合策略等角度分析算法改进方向\custcite{van2007particle}。

若论文中包含关键迭代公式、目标函数或评价指标，可使用公式环境。公式编号按章编排，如\custeqref{eq:update}所示：
\begin{equation}
  v_i^{t+1}=\omega_t v_i^t+c_1 r_1(p_i^t-x_i^t)+c_2 r_2(g^t-x_i^t)
  \label{eq:update}
\end{equation}
其中，$v_i^t$ 表示第 $i$ 个粒子在第 $t$ 次迭代时的速度，$x_i^t$ 表示粒子位置，$p_i^t$ 表示个体历史最优位置，$g^t$ 表示群体历史最优位置，$\omega_t$、$c_1$、$c_2$ 分别表示惯性权重和学习因子，$r_1$、$r_2$ 为随机数。

\section{实验设计与结果分析}

实验设计是算法研究类论文的重要组成部分。论文中应说明测试函数、数据集来源、参数设置、运行次数、评价指标和对比算法。表\ref{tab:datasets}展示了基准测试设置示例。

\begin{table}[H]
  \centering
  \caption{基准测试设置示例}
  \label{tab:datasets}
  \zihao{5}
  \begin{tabular}{cccc}
    \toprule
    测试函数 & 维数 & 搜索空间 & 评价指标 \\
    \midrule
    Sphere & 30 & $[-100,100]^d$ & 最优值、均值、标准差 \\
    Rastrigin & 30 & $[-5.12,5.12]^d$ & 最优值、均值、标准差 \\
    Rosenbrock & 30 & $[-30,30]^d$ & 最优值、均值、标准差 \\
    \bottomrule
  \end{tabular}
\end{table}

实验结果部分可从收敛曲线、最终适应度、稳定性和运行时间等角度展开分析。若结果表格较多，可将完整数据放入附录，正文保留关键对比结果和结论性讨论。根据参考规范，表题置于表上方，图题置于图下方，图表编号按章编排，表格一般采用三线表。
